% GENERATED FILE. Edit canonical lessons, then run npm run generate:books.
% canonical-source-hash: fnv1a64:77d1adfd261f0d72
% canonical-lessons: BN-C97-lobh, BN-C97-pochhondo, BN-C97-shokh, BN-C97-abdar, BN-C97-dabi

\chapter{Greed, Liking, Hobby, Pleading request, Demand}
\label{ch:bn-greed-liking-hobby-pleading-request-demand}
\hlchaptermodality{97}

\begin{chapteropening}
\textbf{By the end of this chapter:} \emph{I can say greed, a liking, a hobby, a pleading request and a demand.}

Complete the last lesson of chapter 97: I can say greed, a liking, a hobby, a pleading request and a demand.
\end{chapteropening}

\section[\texorpdfstring{\bn{লোভ} (lobh)}{লোভ (lobh)}]{\bn{লোভ} (lobh) — greed}

\label{lesson:BN-C97-lobh}



\begin{quote}
Before the new one: say the Bengali for a need, then the Bengali for a necessity.
\end{quote}

\subsection*{You'll want to know: \bn{লোভ}}
\textbf{\bn{লোভ}} — \emph{lobh} — \textquotedblleft{}greed\textquotedblright{}.

\begin{cousinweb}
\textbf{\bn{লোভ}} (\emph{lobh}) is Sanskrit \emph{lobha}, \textquotedblleft{}eager desire, greed\textquotedblright{}, taken over whole as a learned loan, a tatsama.
\end{cousinweb}

\begin{grammarlens}[title={What you've built}]
One more word for this chapter.
\end{grammarlens}

\subsection*{Your turn}
\begin{itemize}
\raggedright
  \item \emph{Say it:} \emph{lobh}
  \item \emph{Say it:} \emph{lobh}, once more
  \item \emph{Say it:} say \emph{lobh}
  \item \emph{Recall:} say the Bengali for a need, then the Bengali for a necessity, then say \emph{lobh} again
\end{itemize}

\begin{tcolorbox}[breakable,colback=teal!4,colframe=teal!35!black,title={Before you move on}]
What does \bn{লোভ} mean? (\textquotedblleft{}Greed\textquotedblright{}.) Say it once more.
\end{tcolorbox}
\section[\texorpdfstring{\bn{পছন্দ} (pôchhôndô)}{পছন্দ (pôchhôndô)}]{\bn{পছন্দ} (pôchhôndô) — a liking, a choice}

\label{lesson:BN-C97-pochhondo}



\begin{quote}
Before the new one: say the Bengali for a necessity, then the Bengali for greed.
\end{quote}

\subsection*{You'll want to know: \bn{পছন্দ}}
\textbf{\bn{পছন্দ}} — \emph{pôchhôndô} — \textquotedblleft{}a liking, a choice\textquotedblright{}.

\begin{cousinweb}
\textbf{\bn{পছন্দ}} (\emph{pôchhôndô}) is a loan, usually traced to Persian \emph{pasand}, \textquotedblleft{}liking, approval\textquotedblright{}.
\end{cousinweb}

\begin{grammarlens}[title={What you've built}]
One more word for this chapter.
\end{grammarlens}

\subsection*{Your turn}
\begin{itemize}
\raggedright
  \item \emph{Say it:} \emph{pôchhôndô}
  \item \emph{Say it:} \emph{pôchhôndô}, once more
  \item \emph{Say it:} say \emph{pôchhôndô}
  \item \emph{Recall:} say the Bengali for a necessity, then the Bengali for greed, then say \emph{pôchhôndô} again
\end{itemize}

\begin{tcolorbox}[breakable,colback=teal!4,colframe=teal!35!black,title={Before you move on}]
What does \bn{পছন্দ} mean? (\textquotedblleft{}A liking, a choice\textquotedblright{}.) Say it once more.
\end{tcolorbox}
\section[\texorpdfstring{\bn{শখ} (shôkh)}{শখ (shôkh)}]{\bn{শখ} (shôkh) — a hobby}

\label{lesson:BN-C97-shokh}



\begin{quote}
Before the new one: say the Bengali for greed, then the Bengali for a liking.
\end{quote}

\subsection*{You'll want to know: \bn{শখ}}
\textbf{\bn{শখ}} — \emph{shôkh} — \textquotedblleft{}a hobby\textquotedblright{}.

\begin{cousinweb}
\textbf{\bn{শখ}} (\emph{shôkh}) is a loan from Arabic \emph{shawq}, \textquotedblleft{}longing, desire\textquotedblright{}, which reached Bengali through Persian. Its Arabic root is \emph{sh-w-q} (Lane). In Bengali it has come to mean \textquotedblleft{}a hobby\textquotedblright{}.
\end{cousinweb}

\begin{grammarlens}[title={What you've built}]
One more word for this chapter.
\end{grammarlens}

\subsection*{Your turn}
\begin{itemize}
\raggedright
  \item \emph{Say it:} \emph{shôkh}
  \item \emph{Say it:} \emph{shôkh}, once more
  \item \emph{Say it:} say \emph{shôkh}
  \item \emph{Recall:} say the Bengali for greed, then the Bengali for a liking, then say \emph{shôkh} again
\end{itemize}

\begin{tcolorbox}[breakable,colback=teal!4,colframe=teal!35!black,title={Before you move on}]
What does \bn{শখ} mean? (\textquotedblleft{}A hobby\textquotedblright{}.) Say it once more.
\end{tcolorbox}
\section[\texorpdfstring{\bn{আবদার} (ābdār)}{আবদার (ābdār)}]{\bn{আবদার} (ābdār) — a pleading request}

\label{lesson:BN-C97-abdar}



\begin{quote}
Before the new one: say the Bengali for a liking, then the Bengali for a hobby.
\end{quote}

\subsection*{You'll want to know: \bn{আবদার}}
\textbf{\bn{আবদার}} — \emph{ābdār} — \textquotedblleft{}a pleading request\textquotedblright{}.

\begin{cousinweb}
The origin of \textbf{\bn{আবদার}} (\emph{ābdār}) is not settled: neither CDIAL nor DEDR has an entry for it, and a search of their word lists by form and by gloss finds nothing close.
\end{cousinweb}

\begin{grammarlens}[title={What you've built}]
One more word for this chapter.
\end{grammarlens}

\subsection*{Your turn}
\begin{itemize}
\raggedright
  \item \emph{Say it:} \emph{ābdār}
  \item \emph{Say it:} \emph{ābdār}, once more
  \item \emph{Say it:} say \emph{ābdār}
  \item \emph{Recall:} say the Bengali for a liking, then the Bengali for a hobby, then say \emph{ābdār} again
\end{itemize}

\begin{tcolorbox}[breakable,colback=teal!4,colframe=teal!35!black,title={Before you move on}]
What does \bn{আবদার} mean? (\textquotedblleft{}A pleading request\textquotedblright{}.) Say it once more.
\end{tcolorbox}
\section[\texorpdfstring{\bn{দাবি} (dābi)}{দাবি (dābi)}]{\bn{দাবি} (dābi) — a demand}

\label{lesson:BN-C97-dabi}



\begin{quote}
Before the new one: say the Bengali for a hobby, then the Bengali for a pleading request.

\emph{Read it:} \textbf{\bn{আমার} \bn{নাম} \bn{রাম।} \bn{আমি} \bn{ভালো।}} without stopping
\end{quote}

\subsection*{You'll want to know: \bn{দাবি}}
\textbf{\bn{দাবি}} — \emph{dābi} — \textquotedblleft{}a demand\textquotedblright{}.

\begin{cousinweb}
\textbf{\bn{দাবি}} (\emph{dābi}) is a loan from Arabic \emph{daʿwā}, \textquotedblleft{}claim, lawsuit\textquotedblright{}, which reached Bengali through Persian. Its Arabic root is \emph{d-ʿ-w} (Lane).
\end{cousinweb}

\begin{grammarlens}[title={What you've built}]
One more word for this chapter.
\end{grammarlens}

\subsection*{Your turn}
\begin{itemize}
\raggedright
  \item \emph{Say it:} \emph{dābi}
  \item \emph{Say it:} \emph{dābi}, once more
  \item \emph{Say it:} say \emph{dābi}
  \item \emph{Recall:} say the Bengali for a hobby, then the Bengali for a pleading request, then say \emph{dābi} again
\end{itemize}

\begin{tcolorbox}[breakable,colback=teal!4,colframe=teal!35!black,title={Before you move on}]
What does \bn{দাবি} mean? (\textquotedblleft{}A demand\textquotedblright{}.) Say it once more.
\end{tcolorbox}
